\label{app:repro}

\section{Repository Structure}
\label{app:structure}

All code lives under \path{code/} in the project repository. The Phase~1 components are:
\begin{lstlisting}
code/
  causal_moe/
    data/          raw.py, clustering.py, mesh.py, windows.py, splits.py,
                   candidate_edges.py (direct / 2hop / full pools)
    causal_graph/  season.py   BSISO index and season mask
                   nodes.py    300 / 596 node series, distances
                   pcmci.py    Stage A (PC1 + MCI, ParCorr, BH-FDR)
                   validate.py Stage B (decoys, lag mask, top-k, persistence)
                   prior.py    Stage C (pcmci / random / corr; restrict / init / regularize)
    splitter/      dirgnn.py (invariance-trained splitter), cia.py, gsina.py
    experts/       expert.py (per-region expert), pool.py (expert blending)
    drift/         rodionov.py (STARS), channels.py, archive.py, forgetting.py
    baselines/     simple.py
  scripts/         phase1_build_graph.py, phase1_run_stage_c.py,
                   phase1_stage_d_drift.py, generate_phase1_figures.py, ...
  tests/           automated tests (including test_pcmci.py)
  results/phase1/  frozen graph, check outputs, Stage C and D results, figures
thesis/            this LaTeX project
\end{lstlisting}

\section{Result Files Behind the Reported Numbers}
\label{app:resultmap}

\begin{small}
\renewcommand{\arraystretch}{1.15}
\begin{longtable}{>{\raggedright\arraybackslash}p{6.0cm}>{\raggedright\arraybackslash}p{8.2cm}}
\caption{Mapping from reported results to result files (paths relative to \texttt{code/results/}).}\label{tab:resultmap}\\
\toprule
\textbf{Result} & \textbf{File} \\
\midrule
\endfirsthead
\toprule
\textbf{Result} & \textbf{File} \\
\midrule
\endhead
\bottomrule
\endfoot
BSISO variance, season statistics (Table~\ref{tab:bsiso}, Fig.~\ref{fig:season}) & \path{phase1/bsiso_season.json}, \path{.npz} \\
Stage A graph (Table~\ref{tab:stagea}) & \path{phase1/stage_a_50.json} \\
B1 decoys (Fig.~\ref{fig:b1}) & \path{phase1/b1_decoy_pass.json} \\
B2, B3, prior, flags (Figs.~\ref{fig:checks}--\ref{fig:b3}) & \path{phase1/causal_graph_50.json} \\
B4 (Fig.~\ref{fig:b4}) & \path{phase1/stage_a_100.json}, \path{causal_graph_50.json} \\
No-prior pools, season (Table~\ref{tab:pool}) & \path{phase1/stage_c/{direct,2hop,full\}\_none.json} \\
No-prior pools, year-round (Table~\ref{tab:pool}) & \path{step7_all_places_direct_lead7.json}, \path{step7_all_places_2hop.json}, \path{step7_all_places_full.json} \\
Stage C, 7 days (Table~\ref{tab:stagec}, Fig.~\ref{fig:stagec7}) & \path{phase1/stage_c/full_{none,pcmci_restrict\}.json} \\
Stage C, 1 day (Tables~\ref{tab:stagec},~\ref{tab:prior-size}, Fig.~\ref{fig:stagec1}) & \path{phase1/stage_c_lead1/full_{none,pcmci_restrict\}.json} \\
Stage D (Table~\ref{tab:staged}, Fig.~\ref{fig:staged}) & \path{phase1/stage_d_drift_w{1,3\}.json} \\
Alternative splitters (Table~\ref{tab:alt-splitters}) & \path{step4_place22.json}, \path{step_cia_place22_direct_w1.0.json}, \path{step_gsina_place22_direct_seed0.json} \\
Drift detection (Table~\ref{tab:p2-drift}) & \path{step5_drift_place{22,12,29\}\_direct.json} \\
Expert lifecycle (Table~\ref{tab:p2-lifecycle}) & \path{step6_lifecycle_place{22,12,29\}\_direct.json} \\
\end{longtable}
\end{small}

The figures of Chapter~\ref{chap:results} were produced by \path{scripts/generate_phase1_figures.py}, except Figure~\ref{fig:pool}, produced by the project's earlier figure script, and Figure~\ref{fig:stagec1}, produced for this thesis by \path{thesis/scripts/make_stage_c_figures.py} from the Stage~C 1-day result files. The paired tests of Section~\ref{sec:res-stagec} were computed from the per-region values in the same files.

\section{Step-by-Step Reproduction Guide}
\label{app:steps}

The following steps reproduce every Phase~1 result, starting from the repository root. Install the pinned dependencies, and Tigramite for the PCMCI verification tests; place the gridded archive under \path{data/raw/rt_preprocessed_archive/} (six \texttt{\{field\}\_rt\_1979\_2022.npz} files and \path{parameters/sst_ocean_mask.npz}); then run, from \path{code/}:
\begin{lstlisting}
pip install -r requirements.txt tigramite
python -m pytest tests                                   # automated tests
python scripts/build_mesh_cache.py                       # 7-day feature cache
python scripts/build_mesh_cache.py --lead-time-days 1 \
    --out cache/windowed_clustered50_lead1_v2.npz        # 1-day feature cache
python scripts/phase1_build_graph.py                     # Stages A and B (~2.5 h)
python scripts/phase1_run_stage_c.py --pools full --priors none,pcmci --modes restrict
python scripts/phase1_run_stage_c.py --pools full --priors none,pcmci --modes restrict \
    --cache-path cache/windowed_clustered50_lead1_v2.npz \
    --out-dir results/phase1/stage_c_lead1               # 1-day lead
python scripts/phase1_run_stage_c.py --workers 12        # full 27-configuration grid
python scripts/phase1_stage_d_drift.py --window-seasons 1
python scripts/phase1_stage_d_drift.py --window-seasons 3
python scripts/generate_phase1_figures.py                # figures from saved results
\end{lstlisting}
The thesis itself is built from \path{thesis/} with \texttt{latexmk -pdf mythesis.tex} (see \path{thesis/README.md}).

\section{Decision Log}
\label{app:decisions}

\begin{small}
\renewcommand{\arraystretch}{1.2}
\begin{longtable}{>{\raggedright\arraybackslash}p{0.7cm}>{\raggedright\arraybackslash}p{3.8cm}>{\raggedright\arraybackslash}p{5.2cm}>{\raggedright\arraybackslash}p{4.5cm}}
\caption{Methodological decisions: the choice made, the alternatives, and open items.}\label{tab:decisions}\\
\toprule
\textbf{ID} & \textbf{Decision} & \textbf{Choice and rationale} & \textbf{Alternatives / status} \\
\midrule
\endfirsthead
\toprule
\textbf{ID} & \textbf{Decision} & \textbf{Choice and rationale} & \textbf{Alternatives / status} \\
\midrule
\endhead
\bottomrule
\endfoot
D1 & Prior granularity & Region level; the splitter selects whole regions and field choice within a region is unstable & Field- or lag-level prior; kept in graph file \\
D2 & CI test & ParCorr; fast, exact for linear-Gaussian, verified & CMI-kNN, GPDC for nonlinear dependence (future) \\
D3 & Source conditioning & Own past for non-target nodes; full parents for target nodes & Full graph for all 300 nodes (higher cost) \\
D4 & FDR scope & Per target over all 18{,}000 candidates & Over survivors only (less conservative); over all targets jointly \\
D5 & Lag range & 1--60 days, lag 0 excluded & $\taumin=L$ for lead-matched graph (E15) \\
D6 & B2 failures & Removed from prior, kept in graph & Diagnostic only; or hard delete \\
D7 & B3 ranking & $|\MCI|$ & Ablation-based importance; bootstrap frequency \\
D8 & B4 failures & Flag, keep in prior & Exclude flagged links (E25) \\
D9 & Region means & Unweighted cell means & Area-weighted ($\cos\varphi$); difference $<$13\% in weights \\
D10 & Season & Activity-based, target days only & Fixed May--Oct; full-year (E20) \\
D11 & Anomaly climatology & Upstream, 1979--2014 & Refit on 1979--2012 (E23) \\
D12 & Forecast type & Direct, separate model per lead & Autoregressive roll-out \\
D13 & First Stage C comparison & Restrict, full pool, 7 and 1 days & Remaining grid cells (E13--E19) \\
D14 & Data provenance & Archive as received & Confirm originating products and licence \\
D15 & Phase~2 expert structure & Open & Regional, per-node, regime or hybrid experts \\
D16 & Phase~2 datasets & IMDAA, ERA5, ISD (planned) & Versions, periods and licences to confirm \\
\end{longtable}
\end{small}
